% Method and experiment design grounded in src/eviseq_new.

\section{Method}
\label{sec:method}

\subsection{Problem setting and design goal}
\label{sec:setting}

Let $x=(x_1,\ldots,x_S)$ be a source document and
$y=(y_1,\ldots,y_T)$ its reference summary.  The proposed system composes a
pretrained source encoder with a pretrained causal decoder whose tokenizers,
hidden widths, and pretraining objectives may differ.  The encoder produces a
final hidden sequence $H^{(L)}\in\mathbb{R}^{S\times d_e}$ and a small set of
intermediate sequences.  The decoder predicts $y_t$ from the causal prefix
$y_{<t}$ and a source memory in its own width $d_d$.

The interface has two requirements.  It must keep every source position that
survives the configured tokenizer and length limit available to cross-attention,
because hard source selection can remove an entity or a number needed later.
It must also provide a controllable retrieval signal without rewriting the
content representation so aggressively that a pretrained vocabulary head or a
copy route loses its lexical anchor.  \EviSeq addresses these requirements with
the \AFMR{} bridge.  The bridge returns one full source tensor for retrieval,
one aligned final-state tensor for values, and one additive source prior.  The
value tensor shares source positions with the retrieval tensor; there is no
pruned or separately selected source bank and no target-derived evidence
representation.

\subsection{\AFMR interface}
\label{sec:afmr-method}

\paragraph{Controller.}
The bridge summarizes the visible source and the fixed decoder instruction
separately.  With $C\in\{0,1\}^S$ the source-content mask and
$M_p\in\{0,1\}^{S_p}$ the prompt mask, we compute
\begin{equation}
  \begin{aligned}
    \bar h&=\MaskedMean(H^{(L)},C),\\
    \bar p&=\MaskedMean(P,M_p).
  \end{aligned}
\end{equation}
where $P$ is obtained by looking up the decoder prompt tokens.  The two means
are independently RMS-normalized, projected to a controller vector
$c\in\mathbb{R}^{d_c}$, and combined with three features available at
inference: the number $N=\sum_i C_i$ of visible content tokens, the requested
output budget $K$, and their ratio.  Let $K_+=\max(K,0)$.  In the base profile
$d_c=256$.  Define the inference-only budget vector
\begin{equation*}
  \phi(N,K)=\left[\begin{array}{c}
    \log(1+N)\\[-1pt]
    \log(1+K_+)\\[-1pt]
    K_+/\max(N,1)
  \end{array}\right].
\end{equation*}
Then
\begin{equation}
  \begin{aligned}
    c&=\RMSNorm\!\big(W_h\RMSNorm(\bar h)\\[-2pt]
      &\quad +W_p\RMSNorm(\bar p)\\[-2pt]
      &\quad +\operatorname{MLP}(\phi(N,K))\big).
  \end{aligned}
  \label{eq:controller}
\end{equation}
The reference summary and its length are never used to construct $c$.

\paragraph{Token-wise depth adaptation.}
The default bridge exposes the final four encoder states.  For source position
$i$ and tap $j$, let $\widetilde H_{i,j}$ be the independently RMS-normalized
state.  A shared token scorer and a controller preference produce
\begin{equation}
  \begin{aligned}
    \pi_{i,j}&=\Softmax_j\!\left(u_d^{\top}\widetilde H_{i,j}+(W_dc)_j\right),\\
    D_i&=\sum_{j=1}^{4}\pi_{i,j}\widetilde H_{i,j}.
  \end{aligned}
  \label{eq:depth}
\end{equation}
The softmax is over depth for each source token, rather than over source
positions.  A low-rank correction is applied around the final-state anchor:
\begin{equation}
  H_i^*=H_i^{(L)}+g_d(c)\,B_d\,\operatorname{SiLU}\!\left(A_d
  [D_i-\widetilde H_{i,4}]\right),
  \label{eq:depth-residual}
\end{equation}
where $A_d$ has the base-profile rank 128, $B_d$ is zero-initialized, and
$g_d(c)=0.15\,\sigma(\cdot)$ is initialized at 0.02.  The zero output factor
makes the initial memory equal to the final-state path while allowing gradients
to reach the residual and, after an update, the depth scorer and router.

\paragraph{Cross-space feature adaptation and value anchor.}
The base map $\Pi$ is identity when $d_e=d_d$ and an orthogonally initialized
trainable projection otherwise.  The retrieval memory is
\begin{equation}
  \begin{aligned}
    Z_i&=\Pi(H_i^*)+G_f(c)\odot B_f\,\operatorname{SiLU}\\[-2pt]
       &\quad\left(A_f\RMSNorm(H_i^*)\right),
  \end{aligned}
  \label{eq:feature}
\end{equation}
where $A_f,B_f$ form the base-profile rank-256 correction and the feature-wise gate
$G_f(c)=0.20\,\sigma(\cdot)$ is initialized at 0.02.  Its output factor is also
zero-initialized.  Independently, the value anchor is
\begin{equation}
  V_i^0=\Pi(H_i^{(L)}).
  \label{eq:value-anchor}
\end{equation}
Thus $Z$ can learn a retrieval-friendly representation while $V^0$ retains the
final encoder content.  This is a local graph constraint, not a claim that the
values are immutable during full fine-tuning: shared encoder and projection
parameters still receive gradients from both paths.

\paragraph{Soft multi-scale source prior.}
\AFMR retains all valid positions and constructs a relative prior from overlapping
content-token windows with base-profile widths $\mathcal W=\{32,128,512\}$ and
50\% overlap.  A task recipe may register a different capacity profile, but
the resolved architecture and its dimensions are part of that run's manifest.
For each window, a masked mean of $Z$ is scored using the controller and a
scale embedding.  Scores are centered within a document and scale, lifted to
tokens by overlap-add, and mixed across available scales.  The focus strength
and temperature are bounded by
\begin{equation}
  \begin{aligned}
    \lambda(c)&=1.0\,\sigma(w_\lambda^\top c+a_\lambda),\\
    \tau(c)&=0.5+1.5\,\sigma(w_\tau^\top c+a_\tau).
  \end{aligned}
  \label{eq:focus-gates}
\end{equation}
The resulting score is $r_i=\lambda(c)\tanh(s_i/\tau(c))$.  On content
positions $\mathcal C$, the source bias is centered by its log-mean-exp:
\begin{equation}
  \begin{aligned}
    b_i&=r_i-\log\left(\frac{1}{|\mathcal C|}
      \sum_{k\in\mathcal C}\exp r_k\right),\quad i\in\mathcal C,\\
    b_i&=0\quad i\notin\mathcal C.
  \end{aligned}
  \label{eq:source-bias}
\end{equation}
Padding is masked separately by cross-attention.  This prior changes relative
source preference but does not prune the memory or determine the final
attention mass by itself.  All reductions in the focus path use FP32 before the
bias is cast to the attention dtype.

\subsection{Decoder source access and grounded copying}
\label{sec:decoder-method}

\paragraph{Copied cross-attention.}
Every causal decoder layer keeps its pretrained self-attention and then applies
a copied cross-attention block.  Query, key, value, output, and normalization
projections are initialized from the corresponding self-attention modules.
For layer $\ell$ and target position $t$,
\begin{equation}
  \begin{aligned}
    h'_{\ell,t}&=h_{\ell,t}+g_\ell\,\mathcal A_{\ell,t},\\
    g_\ell&=\sigma(\theta_\ell),
  \end{aligned}
  \label{eq:cross-attn}
\end{equation}
where $\mathcal A_{\ell,t}$ is multi-head attention with query
$\RMSNorm(h_{\ell,t})$, keys from $Z$, values from $V^0$, additive bias $b$,
and mask $M$.  The gate $g_\ell$ is initialized at 0.10 and bounded by 1.0.
In expanded form, the attention logits contain
$QK(Z)^\top/\sqrt{d_h}+b$ and the padding mask.
The same source prior is supplied to every layer, but each layer has its own
query and output.  During greedy generation, source keys and values are
prepared once per layer and reused through the cache.

\paragraph{Grounded copy and mixture likelihood.}
The copy head uses fast-tokenizer character offsets to align visible source
spans with decoder vocabulary tokens.  It pools contextual source keys,
combines them with lexical decoder embeddings, and computes a decoder-query
attention distribution over eligible aligned source-token rows
$j\in\mathcal J$.  The source prior is pooled with the same offset weights into
$\widetilde b_j$, so
\begin{equation}
  A^{\mathrm{copy}}_{t,j}=\Softmax_{j\in\mathcal J}\!\left(
  \frac{q_t^{\top}k_j}{\sqrt{d_k}}+\widetilde b_j\right).
  \label{eq:copy-attn}
\end{equation}
Repeated source occurrences are marginalized:
\begin{equation}
  P_{\mathrm{copy}}(v\mid t)=\sum_{j\in\mathcal J:\,\operatorname{tok}_j=v}
  A^{\mathrm{copy}}_{t,j}.
  \label{eq:copy}
\end{equation}
The learned gate $\rho_t$ mixes this distribution with the vocabulary head
applied to the decoder hidden state:
\begin{equation}
  \begin{aligned}
    P(y_t=v\mid y_{<t},x)&=\rho_tP_{\mathrm{copy}}(v\mid t)\\
      &\quad +(1-\rho_t)P_{\mathrm{LM}}(v\mid h_t).
  \end{aligned}
  \label{eq:mixture}
\end{equation}
The copy gate is initialized at 0.05.  If no eligible source token is aligned,
the LM logits are used exactly.  The alignment reads only the source and
tokenizers; it does not read the reference summary or evidence labels.

\subsection{Training, gradient flow, and inference parity}
\label{sec:training-method}

The main objective is the token-level negative log likelihood of the mixture:
\begin{equation}
  \mathcal L_{\mathrm{CE}}=-\sum_{t=1}^{T}
  \log P(y_t\mid y_{<t},x).
  \label{eq:objective}
\end{equation}
There is no contrastive, ranking, distillation, pseudo-reference, or
self-improvement objective in the proposed graph.  During interface warm-up,
the pretrained backbones are frozen while the bridge, copied cross-attention,
and grounded copy parameters learn from the same loss.  Full fine-tuning then
updates the encoder and decoder jointly.  Parameters, gradients, and AdamW
states are stored in FP32; CUDA BF16 autocast is used for heavy operations.
The implementation clips the global norm at the configured training threshold
and computes cross-entropy in chunks to avoid retaining a full vocabulary-logit
tensor.

The live gradient route is
\[
\begin{aligned}
\mathcal L_{\mathrm{CE}}&\rightarrow\{P_{\mathrm{LM}},P_{\mathrm{copy}}\}
\rightarrow\{h_t,Z,V^0,b\}\\
&\rightarrow\{\text{decoder},\text{AFMR},\text{encoder}\}.
\end{aligned}
\]
No detach is placed on $Z$, $V^0$, $b$, the depth mixture, or the controller.
The value-anchor separation therefore constrains the direct route from the
retrieval residual to values, but it does not block encoder gradients.

Teacher forcing and inference share the encoder, prompt embeddings, requested
budget, \AFMR bridge, source memory, source prior, and copy alignment.  Official
evaluation uses one beam, no sampling, and autoregressive argmax with
\texttt{temperature=0}, \texttt{top\_k=0}, and \texttt{top\_p=1}; these fields remain available for a
separate candidate-generation path but never alter training.  The encoder and
bridge run once per example, and cached cross-attention states advance with the
generated prefix.

\section{Experiments}
\label{sec:experiments}

\subsection{Research questions and comparisons}
\label{sec:rqs}

The evaluation follows one chain across \PubMed, \ArXiv, \BookSum, and
\GovReport.  RQ1 compares \EviSeq with fine-tuned decoder-only controls on
source support.  RQ2 compares the same systems and \TfiveGemma on overlap and
semantic quality.  RQ3 uses train-from-pretrained ablations with a fixed
decoder: (i) no encoder, implemented as a decoder-only source-prefix control;
(ii) no \AFMR bridge, using the direct final-state projection with zero source
prior; and (iii) no grounded copy, using the complete \AFMR and decoder path but
the vocabulary head alone.  RQ4 holds the decoder and interface fixed while
replacing the anchor encoder \PPLXEmbed with \QwenCausal, \QwenEmbed, and
\NemotronEmbed.  Every RQ includes all four datasets; missing manifests or
failed runs remain explicitly missing rather than being replaced by a score
from another split.

The decoder-only comparison suite comprises \QwenZeroSix, \QwenFour,
\QwenEight, \LlamaThreeTwoThree, \LlamaThreeEight, \NemotronThree, and
\NemotronEight.  Each model is fine-tuned locally on the same source--target
records.  \TfiveGemma is also fine-tuned on each dataset and is treated as an
encoder--decoder reference, not as a prompt-only system.  For an ablation or
encoder-transfer row, all non-target components and optimization budgets are
held fixed; the full baseline suite is reserved for the main RQ1/RQ2 tables.

\subsection{Datasets and preprocessing}
\label{sec:datasets}

\PubMed contains biomedical articles paired with abstracts; \ArXiv contains
scientific papers paired with abstracts; \BookSum contains long narrative
documents with human-written summaries at multiple granularities
\citep{kryscinski2022booksum}; and \GovReport contains long U.S. government
reports with expert-written summaries \citep{huang2021efficient}.  The official split and task variant for each dataset will be
recorded in a manifest.  The current repository has canonical JSONL preparation
for \PubMed and \ArXiv; the \BookSum and \GovReport manifests and field mapping
are pending and therefore cannot support a completed score in this draft.

Each record is normalized to \texttt{id}, \texttt{text}, and \texttt{summary}.  The encoder receives
a fixed task prefix followed by the source and a tokenizer-specific special
sequence.  The content mask excludes the prefix, special tokens, and padding
from depth and focus calculations.  The decoder instruction is fixed per task,
is tokenized once, and is excluded from supervised labels.  Source and target
length limits are selected per dataset and shared by every model in that
dataset; truncation counts are reported.  No reference text is passed to the
model at inference.

\subsection{Metrics and statistical analysis}
\label{sec:metrics}

Primary quality is ROUGE-1, ROUGE-2, and ROUGE-L F1 using one Perl
ROUGE-1.5.5 wrapper and the same preprocessing for all systems
\citep{lin2004rouge}.  Python ROUGE-1.0.0 is retained only for debugging and is
never combined with the primary values.  BERTScore-F1 is the secondary
headline semantic-similarity measure, with precision and recall as diagnostics
\citep{zhang2020bertscore}; the evaluator loads its backbone from a local path
and fails rather than downloading an alternative.

RQ1 uses the native consistency output of the local AlignScore \texttt{nli\_sp}
procedure.  For example $i$, source chunks are paired with each generated
claim sentence, the strongest support is retained for each claim, and the
claim scores are averaged to obtain $C_i=\AlignScore_{\texttt{nli\_sp}}(x_i,y_i)$
\citep{zha-etal-2023-alignscore}.  The dataset estimate is the mean of $C_i$;
higher native consistency is preferred.  For readability we may also display
$H_i=1-C_i$ and its dataset mean as a lower-is-better, direction-transformed
hallucination-oriented value.  This transform is not a probability or a
calibrated percentage of hallucinated content.  AlignScore and BERTScore are
reported with their model/checkpoint identifiers, maximum pair length,
chunking, and sentence splitter.

For every paired comparison, we will compute a 95\% paired bootstrap interval
over identical test IDs and report the number of examples.  A dataset-level
estimate is the mean over its examples; the secondary four-dataset macro
estimate is the unweighted mean of the four dataset-level effects, with a
dataset-aware interval.  Dataset rows are reported first and a macro value is
left pending whenever a dataset is missing.  One primary endpoint is
pre-registered per research question; remaining metrics and model contrasts are
exploratory unless a multiplicity adjustment is specified.  We will also
record generated and reference lengths, length ratio, empty and truncated
outputs, repeated-trigram rate, unique-output rate, source-copy rate, and (when
available) a manually audited factual-error subset.  These diagnostics prevent
an apparent factuality gain from being explained solely by shorter or empty
outputs.

\subsection{Fairness, implementation, and reporting}
\label{sec:fairness}

All main models use the same split IDs, source and target budgets, preprocessing,
seed policy, optimizer family, number of effective training examples, and
greedy decoding constraints within a dataset.  Each model's native source and
decoder input serialization is recorded explicitly; a shared task instruction
does not imply identical tokenization or chat-template behavior.  The AFMR
recipe uses an interface warm-up followed by full fine-tuning; the number of
optimizer updates is matched to the corresponding baseline rather than inferred
from wall-clock time.  The reported checkpoint is \texttt{last.pt}; no test
score selects a best checkpoint.  Batch size and accumulation are recorded per
model and dataset so that the effective batch is explicit.

The experiment manifest will include local model paths or immutable revisions,
tokenizer special IDs, prompts, source/target limits, seeds, parameter counts
(total and trainable), precision, software versions, GPU type, peak VRAM,
training time, generation time, failed runs, and output file checksums.  A score
cell is valid only when its prediction manifest and metric JSON agree on the
dataset split, example IDs, checkpoint, tokenizer, scorer, and decoding
settings.  The four-dataset tables in Section~\ref{sec:results} intentionally
remain blank until these checks pass.
